% Physical unit 046; accepted Session 03 after fresh 400-dpi authority replay.
\GordanUnitStart{046}{Vita}
\begingroup\small
\begin{center}
{\large\bfseries L I F E.}
\end{center}

I was born, Paulus Albertus Gordan, attached to the Evangelical faith, at Breslau on 27 April 1837, my father being David, a merchant, my mother Friderica of the Friedenthal family, whom, taken from me by premature death while I was still a boy, I greatly mourn; I greatly rejoice that my father still lives.

Taught the elements of letters by private teachers until my twelfth year, I then attended the Friedrich Gymnasium at Breslau, under the rectorship of Vimmer, who still flourishes there. Two years later, after I had completed the fourth and partly the third class, I was moved by my father to devote myself to trade, and learned, first in Berlin and then in Breslau in my father's affairs, the things necessary for mercantile business; nevertheless I was always drawn to letters by a certain love. At that time I also attended the disciplines of the mercantile school. After I had completed this course, both in order to learn the French language better and to become acquainted with foreign lands, I was sent by my father to Geneva, where I stayed two years.

Having returned from Geneva to my father's business, since from boyhood I had embraced mathematics with the greatest delight, and following this inclination, I was privately instructed in this discipline by Professor Schellbach. Moved by that most excellent man to devote myself wholly to mathematical studies, I left my father's business and resolved to attend the University of Berlin, although I had not yet obtained the certificate of maturity. After I had diligently heard Professor Kummer there for three semesters, I went to Neisse, where, after spending a semester in the first class of the gymnasium, I was furnished with the certificate of maturity.

During the four years completed since that time I devoted myself to mathematical study; first at Breslau, then at Koenigsberg, and then again at Breslau, I attended these lectures:
\[
\begin{array}{ll}
\text{Mathematics:}& \text{Ill. Kummer, Joachimsthal, Richelot, Rosenhain, Schroeter, Galle.}\\
\text{Physics:}& \text{Ill. Frankenhein, Neumann.}\\
\text{Astronomy:}& \text{Ill. Galle.}\\
\text{Mineralogy:}& \text{Ill. Roemer.}\\
\text{Philosophy:}& \text{Ill. Braniss.}\\
\text{Politics:}& \text{Ill. Tellkampf.}\\
\text{History:}& \text{Ill. Junkmann.}
\end{array}
\]
To all these men, who have deserved very well of me, I have the greatest thanks and shall always have them.

In the literary contest instituted by the most ample philosophical order of Alma Viadrina,
\begin{center}
\textbullet On the geodetic line\textbullet
\end{center}
I entered the arena, and this little work composed by me on that question was crowned with the prize on the third day before the Nones of August of the preceding year.

\GordanOrnament
\endgroup
